\section{4. Funktionen}

\subsection{Konzept}

\Def \key{Funktion}: Input $\to$ 1 Output aus \key{Zielbereich} ($range(f)$)
\begin{align*}
    f: domain(f) & \to range(f) \\
    x & \mapsto y = f(x)
\end{align*}

\Def \key{Definitionsbereich} $\DD(f) = domain(f)$: alle mögliche Inputs

\Def \key{Wertebereich}/\key{Bildmenge} $\WW(f) = image(f)$: realisierte Outputs

\Def input: \key{unabhängige Variable}/\key{Argument}, output: \key{abhängige Variable}

\Def Sei $f: \DD(f) \to \RR$, $\DD(f)\ne \varnothing$, $D' \subset \DD(f)$.
\key{Einschränkung}/\key{Restriktion} von $f$ auf $D'$ ist $f\rvert_{D'} : D' \to \RR$
mit $f\rvert_{D'}(x) = f(x) \forall x \in D'$

\Def Sei $f : \DD(f) \to \RR$ mit $\DD(f)\ne \varnothing$:
\vspace{-0.6em}
\begin{itemize}
    \item $f$ \key{nach oben beschränkt} falls $M \in \RR$ mit $\forall x \in \DD(f) : f(x)\le M$
    \item $f$ \key{nach unten beschränkt} falls $M \in \RR$ mit $\forall x \in \DD(f) : f(x)\ge M$
    \item $f$ \key{beschränkt} falls $M \in \RR$ mit $\forall x \in \DD(f) : |f(x)|\le M$
\end{itemize}

\Def Sei $f : X \to Y$. $f$. $f$ ist:
\vspace{-0.6em}
\begin{itemize}
    \item \key{injektiv}, falls $\forall x_1, x_2 \in X (f(x_1) = f(x_2) \implies x_1 = x_2)$
    \item \key{surjektiv}, falls $\forall y \in Y \, \exists x \in X : f(x) = y$
    \item \key{bijektiv}, falls $f$ injektiv \textit{und} surjektiv ist
\end{itemize}

\Def $f: X \to Y$, $g: Y \to Z$. \key{Verknüpfung} $g \circ f : X \to Z$ mit $\forall x \in X : g \circ f(x) = g(f(x))$

\Def \key{Identität}(\key{-sfunktion}: $\forall x \in X : id_X(x) = x$

\Satz Falls $f: X \to Y$ bijektiv, ex. eindeutiges $g: Y \to X$ mit $g \circ f = id_X$ und $f \circ g = id_Y$.
$g$ heisst \key{Inverse}/\key{Umkehrfunktion}, $f^{-1}$

\begin{minipage}{\linewidth}
\Def $f: \DD(f) \subset \RR \to \RR$ heisst:
\begin{itemize}
    \item \key{monoton wachsend}: $\forall x_1, x_2 \in X (x_1 < x_2 \implies f(x_1)\le f(x_2))$
    \item \key{streng mon. wachs.}: $\forall x_1, x_2 \in X (x_1 < x_2 \implies f(x_1) < f(x_2))$
    \item \key{monoton fallend}: $\forall x_1, x_2 \in X (x_1 < x_2 \implies f(x_1)\ge f(x_2))$
    \item \key{streng mon. fallend}: $\forall x_1, x_2 \in X (x_1 < x_2 \implies f(x_1) > f(x_2))$
    \item \key{monoton}: monoton wachsend oder fallend
    \item \key{streng monoton}: streng monoton wachsend oder fallend
\end{itemize}
\end{minipage}

\Satz Jede streng monotone Funktion ist injektiv

\Def $f: \RR \to \RR$ heisst:
\vspace{-0.6em}
\begin{itemize}
    \item \key{Ungerade}, falls $\forall x \in \RR : f(-x) = -f(x)$
    \item \key{Gerade}, falls $\forall x \in \RR : f(-x) = f(x)$
\end{itemize}

\Def \key{Graph} der Funktion ist Menge alle Punkte der Form $(x, f(x))$

\Def Graph ist \key{linksgekrümmt}/\key{konvex} falls man beim durchlaufen von L nach R links abbiegt.
Falls rechts \key{rechtsgekrümmt}/\key{konkav}

\subsection{Grenzwerte}

\Def $f: \DD(f) \to \RR$, $x_0 \in \RR$, $\DD(f) \cap (x_0 - \delta, x_0 + \delta)\ne \varnothing \forall \delta > 0$.
Dann $L \in \RR$ \key{Grenzwert}/\key{Limes} von $f(x)$ an Stelle $x_0$, falls
$$\forall \varepsilon > 0 \, \exists \delta > 0 : x \in \DD(f) \cap (x_0 - \delta, x_0 + \delta) \implies |f(x) - L| < \varepsilon$$
$L = \lim\limits_{\substack{x\ge x_0 \\ x \to x_0}} f(x)$ ist \key{rechtsseitigen Grenzwert} falls
$$\forall \varepsilon > 0 \, \exists \delta > 0 : x \in \DD(f) \cap [x_0, x_0 + \delta) \implies |f(x) - L| < \varepsilon$$
$L = \lim\limits_{\substack{x\le x_0 \\ x \to x_0}} f(x)$ ist \key{linksseitigen Grenzwert} falls
$$\forall \varepsilon > 0 \, \exists \delta > 0 : x \in \DD(f) \cap (x_0 - \delta, x_0] \implies |f(x) - L| < \varepsilon$$
Man kann ähnlich $\lim\limits_{\substack{x > x_0 \\ x \to x_0}}$ und $\lim\limits_{\substack{x < x_0 \\ x \to x_0}}$
definieren (mit $()$ anstatt $[]$)
$f$ hat \key{für $x$ gegen $\infty$} den Grenzwert $L = \lim\limits_{x\to\infty}f(x)$, falls $\forall R > 0 : \DD(f) \cap (R, \infty) \ne \varnothing$ und
$$\forall \varepsilon > 0 \, \exists M > 0 : x \in \DD(f) \cap (M, \infty) \implies |f(x) - L| < \varepsilon$$
$f$ hat \key{für $x$ gegen $-\infty$} den Grenzwert $L = \lim\limits_{x\to\infty}f(x)$, falls $\forall R > 0 : \DD(f) \cap (-\infty, -R) \ne \varnothing$ und
$$\forall \varepsilon > 0 \, \exists M > 0 : x \in \DD(f) \cap (-\infty, -M) \implies |f(x) - L| < \varepsilon$$
$f$ hat in $x_0$ den \key{uneigentlichen Grenzwert} $\lim\limits_{x \to x_0}f(x) = \infty$ falls
$$\forall N > 0 \, \exists \delta > 0 : x \in \DD(f) \cap (x_0 - \delta, x_0 + \delta) \implies f(x) \ge N$$
$f$ hat in $x_0$ den \key{uneigentlichen Grenzwert} $\lim\limits_{x \to x_0}f(x) = -\infty$ falls
$$\forall N > 0 \, \exists \delta > 0 : x \in \DD(f) \cap (x_0 - \delta, x_0 + \delta) \implies f(x) \le -N$$

\Satz $f: \DD(f) \to \RR$. $\lim\limits_{x \to x_0} f(x) = L$ $\iff$ für jede konv. Folge $(x_n)_{n \in \NN_0}$,
die gegen $x_0$ konvergiert, gilt $\lim\limits_{n\to\infty} f(x_n) = L$

\Satz \key{Rechenregeln}: $\lim\limits_{x\to c} f(x) = K (\ne \pm \infty)$, $\lim\limits_{x\to c}g(x) = L (\ne \pm \infty)$, $A \in \RR$:
\vspace{-0.6em}
\begin{itemize}
    \item $\lim\limits_{x\to c}(f(x) + g(x)) = K + L$
    \item $\lim\limits_{x\to c}(f(x) - g(x)) = K - L$
    \item $\lim\limits_{x\to c}(f(x) \cdot g(x)) = K \cdot L$
    \item $\lim\limits_{x\to c}(A \cdot f(x)) = A \cdot K$
    \item Falls $L \ne 0$: $\lim\limits_{x\to c}\frac{f(x)}{g(x)} = \frac{K}{L}$
\end{itemize}

\subsection{Elementare Funktionen}

\Def \key{Lineare Funkt.}: $y = f(x) = mx + q$, $x \mapsto mx + q$. Eigenschaft, \key{$\WW$ durch $\DD$ konstant}:
$\frac{f(x+y)-f(x)}{x+y-x} = \frac{mx + my + q -mx - q}{y} = \frac{my}{y} = m$

\Def \key{Exponentialfunktion}: $\exp(x) = e^x = \lim\limits_{n\to\infty} \left(1 + \frac{x}{n}\right)^n$

\Satz $\exp: \RR \to \RR$ ist bijektiv, streng monoton wachsend, stetig und
\vspace{-0.6em}
\begin{itemize}
    \item $\exp(0) = 1$
    \item $\exp(-x) = \exp(x)^{-1} \, \forall x \in \RR$
    \item $\exp(x + y) = \exp(x)\cdot\exp(y) \, \forall x,y \in \RR$
\end{itemize}

\Def Exponentialfunktion mit \key{Anfangswert} $z_0$ und \key{Wachstumfaktor}/\key{Basis} $a$: $z = f(t) = z_0 \cdot a^t = z_0 \cdot e^{kt}$
Eigenschaft: \key{relative Zunahme konstant}: $\frac{z(t + s)}{z(t)} = \frac{z_0 \cdot a^{t + s}}{z_0 \cdot a^t} = a^s$

\Def Eindeutige Inverse zu $\exp$ heisst (\key{natürlicher}) \key{Logarithmus}: $\log(x): \RR^+ \to \RR$

\Satz Logarithmus $\log: \RR^+ \to \RR$ ist bijektiv, streng monoton wachsend, stetig und:
\vspace{-0.6em}
\begin{itemize}
    \item $\log(1) = 0$
    \item $\log(x^{-1}) = -\log(x) \, \forall x \in \RR^+$
    \item $\log(x \cdot y) = \log(x) + \log(y) \, \forall x,y \in \RR^+$
\end{itemize}

\Def $g(x) = s(u(x)) = s \circ u(x)$: \key{zusammengesetzte Funktion}. $s$: \key{äussere Funk.}. $u$: \key{innere Funk.}

\Satz Für trigonometrische Funktionen gilt:
\vspace{-0.6em}
\begin{itemize}
    \item $\sin$, $\cos$ und $\tan$ sind \key{periodische} Funktionen
    \item $\tan$ hat für $\omega = \pm \frac{(2n+1)\pi}{2}$, $n \in \NN_0$ eine \key{vertikale Asymptote}
\end{itemize}
(trig. Identitäten $\to$ Abschnitt 10)

\Def $f(t) = A \sin(B(t+h)) + C$ oder $f(t) = A\cos(B(t+h))+C$ sind \key{Sinusschwingungen}/\key{Sinuswellen} und:
\vspace{-0.6em}
\begin{itemize}
    \item $|A|$ heisst \key{Amplitude}
    \item $|B|^{-1} \cdot 2\pi$ nennen wir \key{Periode}
    \item $|h|$ ist die \key{Phasenverschiebung}
    \item $|C|$ wird \key{Mittelwert} genannt
\end{itemize}

\Def $y = f(x) = k\cdot x^p$, $k,p \in \RR$, heisst \key{Potenzfunktion}

\Def $y = f(x) = \sum_{s=0}^{n} a_s x^s$, $n \in \NN$ und alle $a_s \in \RR$, heisst \key{Polynom}.
$a_s$ heissen \key{Koeffizienten}, $n$ ist der \key{Grad} des Polynoms.

\Def $y = f(x) = \frac{p(x)}{q(x)}$, $p(x)$ und $q(x)$ Polynome, heisst \key{rationale F.}

\subsection{Stetigkeit}

\Def $f: \DD(f) \to \RR$ \key{an der Stelle $x_0$ stetig} wenn:
$\forall \varepsilon > 0 \, \exists \delta > 0 \, \forall x \in \DD (|x - x_0| < \delta \implies |f(x) - f(x_0)| < \varepsilon)$.
$f$ \key{stetig} falls $\forall x_0 \in \DD$ stetig

\Def $f: \DD(f) \to \RR$, $x_0 \in \DD(f)$ ist \key{rechtsstetig} falls
$\lim\limits_{\substack{x\ge x_0 \\ x \to x_0}} f(x) = f(x_0)$.
\key{Linksstetigkeit} ist analog definiert

\Satz $f: \DD(f) \to \RR$, $x_0 \in \DD(F)$. $f$ an Stelle $x_0$ stetig $\iff$
für jede Folge $(x_n)_{n\in\NN_0}$ mit $x_n \to x_0$ konvergiert $(f(x_n))_{n\in\NN_0}$ gegen $f(x_0)$

\Satz \key{Rechenregeln}: für $f,g: D \subset \RR \to \RR$ stetig, $c \in \RR$ gilt
\vspace{-0.6em}
\begin{itemize}
    \item $f + g$ ist stetig
    \item $f - g$ ist stetig
    \item $c \cdot f$ ist stetig
    \item $f \cdot g$ ist stetig
    \item $\frac{f}{g}$ ist stetig, sofern $g \ne 0$
    \item Falls $f: I \subset \RR \to image(f)$, $g: image(f) \to image(g)$ ist
    $g \circ f$ stetig und $\lim\limits_{x\to x_0} g(f(x)) = g(\lim\limits_{x \to x_0}f(x))$
\end{itemize}

\Satz \key{Stetigkeit der Umkehrfunktion}: $I, J \subset \RR$ und $f: I \to J$ stetig und
bijektiv. $f^{-1}: J \to I$ ist auch stetig

\Satz \key{Zwischenwertsatz}: $f: [a, b] \to \RR$ stetig, $c \in [f(a), f(b)]$.
Es gibt ein $x \in [a, b]$ mit $f(x) = c$

\Kor $f: I \subset \RR \to \RR$ stetig, $a, b \in I$ mit $f(a) \le 0$ und $f(b) \ge 0$.
$f$ hat mindestens eine Nullstelle zwischen $a$ und $b$

\Def ein beschränktes und abgeschlossenes Intervall ist \key{kompakt}

\Satz $[a, b]$ kompakt, $(x_n)_{n\in\NN_0}$ in $[a, b]$.
Dann existiert eine Teilfolge $(x_{n_k})_{k\in\NN_0}$ mit $\lim\limits_{k\to\infty} x_{n_k} = x_0$ mit $x_0 \in [a, b]$

\Def $f: D \subset \RR \to \RR$, $x_0 \in D$. Dann liegt an der Stelle $x_0$ ein
\vspace{-0.6em}
\begin{itemize}
    \item \key{lokales Maximum} falls $f(x_0) \ge f(x)$ für alle $x \in D$
    \item \key{lokales Minimum} falls $f(x_0) \le f(x)$ für alle $x \in D$
    \item \key{lokales Extremum} falls lokales Minimum oder Maximum
\end{itemize}

\Satz: \key{Stetige Funktion auf kompaktem Intervall}: $f: I \subset \RR \to \RR$ stetig,
$I$ kompakt. $f$ beschränkt und $f$ nimmt max und min auf $I$ an,
d.h. es gibt $x_{min}, x_{max} \in I$ mit $\forall x \in I : f(x_{min}) \le f(x) \le f(x_{max})$
